\subsection{Five-dimensional pseudo-experiments with both background treatments}

The five-dimensional validation (\S\S\ref{sec:val-fiveaxis}--\ref{sec:val-departure}) reports end-to-end tests of both background treatments. They cover the purity bias in the highest-$W$ column at nominal truth, its reduction with refined negative-weight subtraction, the data difference, and the estimator with higher refinement capacity. That estimator is the \estMultiD{}, which gives the 3D and 4D central values and the sixteen-cell result (\S\ref{sec:fived16}) but not the \resFiveAxis{}; the quoted 2D central value uses the same refinement capacity, in the \estTwoD{}.

The closure of the five-dimensional pseudo-experiments with backgrounds, under purity-based and refined negative-weight subtraction (\S\ref{sec:val-negweight}), is shown cell by cell in Fig.~\ref{fig:bkg-residual-note}.

% Targeted round (DECISION-20261103-joseph-retire-density-gate.md item 3b): the methods paper's Fig. 2 and its caption, verbatim, moved here when the paper's background-closure paragraph (Sec. IV D) moved to the note. pointed to from sec:val-negweight. The same figure file the paper printed (figures/paper_fig_bkg_residual.pdf, unchanged).
\begin{figure}[tbp]
\centering
\includegraphics[width=\linewidth]{paper_fig_bkg_residual}
\caption[Closure residual with backgrounds: purity and negative-weight subtraction]{Development check: purity-based subtraction leaves a bias in the highest-$W$ column at nominal truth, and refined negative-weight subtraction reduces it there, but residuals in some joint cells fail the nominal closure criterion (\S\ref{sec:val-negweight}). Bars: standard error of the mean. The plotted quantity is the nominal-truth fractional closure residual, (unfolded$-$truth)/truth, with backgrounds and deterministic LightGBM. Circles: purity subtraction (9 pseudo-experiments). Squares: unbinned negative-weight subtraction, 100-tree/8-leaf refinement (60). The pseudo-experiments are unpaired. (a)~42 $(\Eavail,W)$ cells grouped by $\Eavail$ bin. (b)~\jtNcells{} coarse joint cells indexed in the $3^{5}$ grid, followed by the integrated total.}
\label{fig:bkg-residual-note}
\end{figure}

\paragraph{Purity bias.} Purity bias is 8--12 statistical standard deviations. Signal-only experiments close. Subtraction is exact per reconstructed bin, but the widest bin ($W\ge\SI{3}{GeV}$, background \SI{13}{\percent} of signal) contains background shape resolved by the unbinned classifier and missed by bin-uniform purity. The within-bin shape term of \S\ref{sec:negweight-estimator} is therefore measurable, and it is direct evidence for revisiting the background treatment of the \resFiveAxis{} (\S\ref{sec:negweight-footing}). The difference between five-dimensional purity and refined negative-weight subtraction on data is the least computationally expensive estimate of its size on data.
% Source: docs/orchestration/OUTCOME-20260925-s5c-purity-background-bias-at-high-W.md; KNOWN_ISSUES 75.
% Heading date (2026-09-25) and the numbers repeated in sec_validation.tex removed (lane 3 App. B #7):
% highest-$W$ column \SIrange{3.5}{4.1}{\percent} low, highest-$\Eavail$, highest-$W$ cell \SI{1.7}{\percent} high.
